\chapter{Parameters and Optimization}
\label{ch:parameters}

This chapter explores the parameter architecture that transforms a theoretical biological model into a calibratable computational system. Parameter design in agent-based biological models faces a fundamental challenge: balancing biological interpretability with optimization tractability while maintaining model identifiability. Our approach embodies a hierarchical parameter philosophy that separates concerns across biological, computational, and optimization scales, enabling systematic model calibration without sacrificing biological insight.

% ======================================================================
\section{Parameter Design Philosophy}
\label{sec:param-philosophy}

The 92-parameter space represents a carefully designed hierarchy that separates biological interpretability from computational optimization concerns. This separation reflects a core principle in biological modeling: parameters should have clear biological meaning when possible, but optimization requires mathematical flexibility that may obscure biological interpretation. Our design resolves this tension through a three-tier architecture that isolates patient-specific clinical variables, system-level model configuration, and process-level behavioral weights.

The parameter hierarchy embodies several design principles that emerged from the constraints of biological modeling and Bayesian optimization. \textbf{Separability}: patient parameters remain fixed for individual simulations while weight parameters vary during optimization, preventing biological confounding of optimization landscapes. \textbf{Dimensionality}: parameters are grouped by their biological role rather than their mathematical properties, enabling domain experts to reason about parameter interactions within familiar biological contexts. \textbf{Identifiability}: each parameter governs a distinct biological process or computational behavior, reducing the risk of parameter redundancy that can destabilize optimization procedures.

This design philosophy ultimately serves the broader goal of model validation: parameters must be interpretable enough for biologists to assess their reasonableness while being mathematically structured enough for optimization algorithms to explore efficiently.

% ======================================================================
\section{Parameter Categories}
\label{sec:param-categories}

The model's parameter space is partitioned into three \texttt{dataclass}-based parameter objects, each inheriting from a common \texttt{Parameters} base class that provides serialization, XML loading, and keyword-argument initialization. This organizational structure reflects the distinct roles these parameter categories play in model configuration and optimization.

\subsection{ModelParameters}

\texttt{ModelParameters} controls the simulation geometry and execution. It contains 4 parameters:

\begin{table}[H]
\centering
\caption{Model parameters.}
\label{tab:model-params}
\begin{tabular}{@{}llp{7cm}@{}}
\toprule
\textbf{Parameter} & \textbf{Default} & \textbf{Description} \\
\midrule
\texttt{volume} & 0.0001 & Simulated tissue volume in mL. Determines the 3D grid dimension via \texttt{get\_dimension\_from\_volume()}. \\
\texttt{block\_size} & 10 & Spatial resolution in $\mu$m per grid cell. \\
\texttt{random\_seed} & 1 & Seed for the \texttt{random.Random} instance; added to MPI rank for per-rank determinism. \\
\texttt{max\_steps} & 500 & Maximum number of simulation ticks before termination. \\
\bottomrule
\end{tabular}
\end{table}

\subsection{PatientParameters}

\texttt{PatientParameters} encodes patient-specific clinical and biological features. It contains 18 parameters:

\begin{table}[H]
\centering
\caption{Patient parameters.}
\label{tab:patient-params}
\begin{tabular}{@{}llp{6.5cm}@{}}
\toprule
\textbf{Parameter} & \textbf{Default} & \textbf{Description} \\
\midrule
\texttt{sex} & \texttt{'F'} & Patient sex (\texttt{'F'} or \texttt{'M'}). Controls hormone chances, CD8+ drift direction, and initial DNA mutations. \\
\texttt{BMI} & 22.0 & Body mass index (kg/m$^2$). Modulates Treg differentiation, macrophage mutation, and NK cell kill rate via BMI weight parameters. \\
\texttt{treatment} & \texttt{'ICI+TKI'} & Treatment modality: \texttt{'None'}, \texttt{'ICI'}, \texttt{'TKI'}, or \texttt{'ICI+TKI'}. \\
\texttt{treatment\_start} & 100 & Simulation step at which treatment begins. \\
\midrule
\multicolumn{3}{@{}l}{\textit{Immune cell concentrations (cells/mL):}} \\
\texttt{ctc\_concentration} & 80\,000 & Cytotoxic T cell (CD8+) concentration. \\
\texttt{neutrophil\_concentration} & 40\,000 & Neutrophil concentration. \\
\texttt{mast\_cell\_concentration} & 40\,000 & Mast cell concentration. \\
\texttt{treg\_concentration} & 40\,000 & Regulatory T cell concentration. \\
\texttt{pdc\_concentration} & 40\,000 & Plasmacytoid dendritic cell concentration. \\
\texttt{th1\_concentration} & 40\,000 & CD4+ Th1 helper cell concentration. \\
\texttt{th2\_concentration} & 40\,000 & CD4+ Th2 helper cell concentration. \\
\texttt{dc\_concentration} & 80\,000 & Dendritic cell concentration. \\
\texttt{m1\_concentration} & 40\,000 & M1 macrophage concentration. \\
\texttt{m2\_concentration} & 40\,000 & M2 macrophage concentration. \\
\texttt{nkl\_concentration} & 160\,000 & Natural killer cell concentration. \\
\texttt{cd4\_concentration} & 80\,000 & Na\"ive CD4+ T cell concentration. \\
\texttt{cd8\_concentration} & 80\,000 & Na\"ive CD8+ T cell concentration. \\
\texttt{NKL\_delta} & 16\,000 & NK cell count difference (used in sex-specific adjustments). \\
\bottomrule
\end{tabular}
\end{table}

\subsection{WeightParameters}

\texttt{WeightParameters} is the largest parameter class, containing all behavioral weights, biases, and system-level tuning knobs. In total, it comprises 84 original parameters plus 8 new glucose parameters, for a total of 92 tunable values. These are detailed in Sections~\ref{sec:weight-params} and~\ref{sec:glucose-params}.

% ======================================================================
\section{Weight Parameters}
\label{sec:weight-params}

Weight parameters control the probability thresholds, scaling factors, and biases that govern agent behaviors. They are organized by the agent type or subsystem they influence. Bias parameters (prefixed with \texttt{b\_}) are additive offsets to the corresponding weight's contribution, allowing the optimization to shift decision boundaries.

\subsection{BMI Weights (4 parameters)}

These weights modulate the influence of body mass index on immune cell behavior:

\begin{table}[H]
\centering
\caption{BMI weight parameters.}
\label{tab:bmi-weights}
\begin{tabular}{@{}llp{6cm}@{}}
\toprule
\textbf{Parameter} & \textbf{Default} & \textbf{Description} \\
\midrule
\texttt{w\_BMI\_on\_treg\_diff} & 0.01 & BMI influence on Treg differentiation probability. \\
\texttt{w\_BMI\_on\_m1\_mutation} & 0.01 & BMI influence on M1 macrophage mutation threshold. \\
\texttt{w\_BMI\_on\_m2\_mutation} & 0.01 & BMI influence on M2 macrophage mutation threshold. \\
\texttt{w\_BMI\_nkl\_kill\_rate} & 0.01 & BMI influence on NK cell kill rate. \\
\bottomrule
\end{tabular}
\end{table}

\subsection{Macrophage M1 Weights (6 weights + 2 biases)}

\begin{table}[H]
\centering
\caption{Macrophage M1 weight parameters.}
\label{tab:m1-weights}
\begin{tabular}{@{}llp{6cm}@{}}
\toprule
\textbf{Parameter} & \textbf{Default} & \textbf{Description} \\
\midrule
\texttt{w\_m1\_mutation} & 1.0 & M1 macrophage mutation threshold scaling. \\
\texttt{b\_m1\_mutation} & 0.0 & Bias for M1 mutation threshold. \\
\texttt{w\_m1\_move} & 1.0 & M1 macrophage look-up movement radius scaling. \\
\texttt{w\_m1\_phagocytosis} & 1.0 & M1 phagocytosis chance scaling. \\
\texttt{w\_m1\_digest} & 1.0 & M1 digestion rate scaling. \\
\texttt{w\_m1\_t\_kill\_rate} & 1.0 & M1 effect on T cell kill rate. \\
\texttt{w\_m1\_th1\_proliferation} & 1.0 & M1 effect on Th1 proliferation. \\
\texttt{w\_th1\_proliferation} & 1.0 & Th1 proliferation scaling (via M1 influence). \\
\texttt{b\_th1\_proliferation} & 0.0 & Bias for Th1 proliferation. \\
\bottomrule
\end{tabular}
\end{table}

\subsection{Macrophage M2 Weights (5 weights + 1 bias)}

\begin{table}[H]
\centering
\caption{Macrophage M2 weight parameters.}
\label{tab:m2-weights}
\begin{tabular}{@{}llp{6cm}@{}}
\toprule
\textbf{Parameter} & \textbf{Default} & \textbf{Description} \\
\midrule
\texttt{w\_m2\_mutation} & 1.0 & M2 macrophage mutation threshold scaling. \\
\texttt{b\_m2\_mutation} & 0.0 & Bias for M2 mutation threshold. \\
\texttt{w\_m2\_move} & 1.0 & M2 macrophage look-up movement radius. \\
\texttt{w\_m2\_t\_kill\_rate} & 1.0 & M2 effect on T cell kill rate. \\
\texttt{w\_m2\_tumour\_growth} & 1.0 & M2 effect on tumour growth promotion. \\
\texttt{w\_m2\_angiogenesis} & 1.0 & M2 effect on angiogenesis. \\
\bottomrule
\end{tabular}
\end{table}

\subsection{Tumor Cell Weights (8 weights + 3 biases)}

\begin{table}[H]
\centering
\caption{Tumor cell weight parameters.}
\label{tab:tumor-weights}
\begin{tabular}{@{}llp{5.5cm}@{}}
\toprule
\textbf{Parameter} & \textbf{Default} & \textbf{Description} \\
\midrule
\texttt{w\_tumor\_apoptosis\_eff} & 1.0 & Effect-driven apoptosis probability scaling. \\
\texttt{w\_tumor\_apoptosis\_dna} & 1.0 & DNA-driven extra suppression chance scaling. \\
\texttt{b\_tumor\_apoptosis} & 0.0 & Bias for apoptosis probability. \\
\texttt{w\_tumor\_growth\_eff} & 0.6 & Effect-driven growth probability scaling. \\
\texttt{w\_tumor\_growth\_dna} & 1.0 & DNA-driven extra proliferation chance scaling. \\
\texttt{b\_tumor\_growth} & 0.0 & Bias for growth probability. \\
\texttt{w\_tumor\_angiogenesis} & 1.0 & Angiogenesis initiation probability scaling. \\
\texttt{b\_tumor\_angiogenesis} & 0.0 & Bias for angiogenesis probability. \\
\texttt{w\_antigen\_presentation} & 1.0 & Antigen presentation chance scaling. \\
\texttt{b\_antigen\_presentation} & 0.0 & Bias for antigen presentation. \\
\texttt{w\_angiogenesis\_tumor\_growth} & 1.0 & Angiogenesis contribution to growth (when blood access is present). \\
\texttt{w\_gene\_pd1\_inhibition} & 1.0 & PD-1/PD-L1 checkpoint pathway inhibition scaling. \\
\bottomrule
\end{tabular}
\end{table}

\subsection{CD4+ T Cell Weights (9 weights + 3 biases)}

\begin{table}[H]
\centering
\caption{CD4+ T cell weight parameters.}
\label{tab:cd4-weights}
\begin{tabular}{@{}llp{5.5cm}@{}}
\toprule
\textbf{Parameter} & \textbf{Default} & \textbf{Description} \\
\midrule
\texttt{w\_cd4\_treg\_diff\_effect} & 1.0 & Effect-driven Treg differentiation scaling. \\
\texttt{w\_cd4\_treg\_diff\_horm} & 1.0 & Hormone-driven Treg differentiation scaling. \\
\texttt{b\_cd4\_treg\_diff} & 0.0 & Bias for Treg differentiation. \\
\texttt{w\_cd4\_th1\_proliferation\_effect} & 1.0 & Effect-driven Th1 proliferation scaling. \\
\texttt{w\_cd4\_th1\_proliferation\_horm} & 1.0 & Hormone-driven Th1 proliferation scaling. \\
\texttt{b\_cd4\_th1\_proliferation} & 0.0 & Bias for Th1 proliferation. \\
\texttt{w\_cd4\_th1\_spawn\_m1} & 1.0 & Th1 spawning of M1 macrophages. \\
\texttt{w\_cd4\_th1\_spawn\_dc} & 1.0 & Th1 spawning of dendritic cells. \\
\texttt{w\_cd4\_th2\_spawn\_m1} & 1.0 & Th2 spawning of M1 macrophages. \\
\texttt{w\_cd4\_th2\_spawn\_t} & 1.0 & Th2 spawning of T cells. \\
\texttt{w\_cd4\_th2\_proliferation\_effect} & 1.0 & Effect-driven Th2 proliferation scaling. \\
\texttt{w\_cd4\_th2\_proliferation\_horm} & 1.0 & Hormone-driven Th2 proliferation scaling. \\
\texttt{b\_cd4\_th2\_proliferation} & 0.0 & Bias for Th2 proliferation. \\
\bottomrule
\end{tabular}
\end{table}

\subsection{Cytotoxic T Cell Weights (5 weights + 2 biases)}

\begin{table}[H]
\centering
\caption{Cytotoxic T cell (CD8+) weight parameters.}
\label{tab:ctc-weights}
\begin{tabular}{@{}llp{5.5cm}@{}}
\toprule
\textbf{Parameter} & \textbf{Default} & \textbf{Description} \\
\midrule
\texttt{w\_cytotoxic\_proliferation} & 1.0 & Proliferation probability scaling. \\
\texttt{w\_cytotoxic\_apoptosis} & 1.0 & Apoptosis probability scaling. \\
\texttt{b\_cytotoxic\_apoptosis} & 0.0 & Bias for cytotoxic T cell apoptosis. \\
\texttt{w\_cytotoxic\_kill} & 1.0 & Tumor cell kill probability scaling. \\
\texttt{b\_cytotoxic\_kill} & 0.0 & Bias for cytotoxic kill probability. \\
\texttt{w\_cytotoxic\_pd1\_inhibition} & 1.0 & PD-1 inhibition susceptibility scaling. \\
\texttt{w\_cytotoxic\_move} & 1.0 & Movement radius scaling for cytotoxic T cells. \\
\bottomrule
\end{tabular}
\end{table}

\subsection{Mast Cell Weights (6 parameters)}

\begin{table}[H]
\centering
\caption{Mast cell weight parameters.}
\label{tab:mast-weights}
\begin{tabular}{@{}llp{5.5cm}@{}}
\toprule
\textbf{Parameter} & \textbf{Default} & \textbf{Description} \\
\midrule
\texttt{w\_mast\_cell\_angiogenesis} & 1.0 & Mast cell contribution to angiogenesis effect. \\
\texttt{w\_mast\_cell\_m1\_mutation} & 1.0 & Mast cell effect on M1 mutation. \\
\texttt{w\_mast\_cell\_t\_kill\_rate} & 1.0 & Mast cell effect on T cell kill rate. \\
\texttt{w\_mast\_cell\_tumour\_apoptosis} & 1.0 & Mast cell contribution to tumour apoptosis. \\
\texttt{w\_mast\_cell\_tumour\_growth} & 1.0 & Mast cell contribution to tumour growth. \\
\texttt{w\_mast\_cell\_spawn\_dc} & 1.0 & Mast cell probability of spawning dendritic cells. \\
\bottomrule
\end{tabular}
\end{table}

\subsection{NK Cell Weights (3 weights + 1 bias)}

\begin{table}[H]
\centering
\caption{Natural killer cell weight parameters.}
\label{tab:nk-weights}
\begin{tabular}{@{}llp{5.5cm}@{}}
\toprule
\textbf{Parameter} & \textbf{Default} & \textbf{Description} \\
\midrule
\texttt{w\_natural\_killer\_kill\_rate} & 1.0 & NK cell direct kill rate scaling. \\
\texttt{b\_natural\_killer\_kill\_rate} & 0.0 & Bias for NK cell kill rate. \\
\texttt{w\_nkl\_t\_kill\_rate} & 1.0 & NK cell effect on T cell kill rate amplification. \\
\bottomrule
\end{tabular}
\end{table}

\subsection{Plasmacytoid DC Weights (6 weights + 1 bias)}

\begin{table}[H]
\centering
\caption{Plasmacytoid dendritic cell weight parameters.}
\label{tab:pdc-weights}
\begin{tabular}{@{}llp{5.5cm}@{}}
\toprule
\textbf{Parameter} & \textbf{Default} & \textbf{Description} \\
\midrule
\texttt{w\_pdc\_nkl\_spawn} & 1.0 & pDC probability of spawning NK cells. \\
\texttt{b\_pdc\_nkl\_spawn} & 0.0 & Bias for pDC NK cell spawning. \\
\texttt{w\_pdc\_angiogenesis} & 1.0 & pDC contribution to angiogenesis effect. \\
\texttt{w\_pdc\_treg\_diff} & 1.0 & pDC effect on Treg differentiation. \\
\texttt{w\_pdc\_t\_proliferation} & 1.0 & pDC effect on T cell proliferation. \\
\texttt{w\_pdc\_t\_kill} & 1.0 & pDC effect on T cell kill rate. \\
\texttt{w\_pdc\_nkl\_kill} & 1.0 & pDC effect on NK cell kill rate. \\
\bottomrule
\end{tabular}
\end{table}

\subsection{Treg Weights (5 parameters)}

\begin{table}[H]
\centering
\caption{Regulatory T cell weight parameters.}
\label{tab:treg-weights}
\begin{tabular}{@{}llp{5.5cm}@{}}
\toprule
\textbf{Parameter} & \textbf{Default} & \textbf{Description} \\
\midrule
\texttt{w\_treg\_t\_kill\_rate} & 1.0 & Treg suppression of T cell kill rate. \\
\texttt{w\_treg\_t\_proliferation} & 1.0 & Treg suppression of T cell proliferation. \\
\texttt{w\_treg\_t\_apoptosis} & 1.0 & Treg induction of T cell apoptosis. \\
\texttt{w\_treg\_activation} & 1.0 & Treg activation threshold scaling. \\
\texttt{w\_treg\_dc\_phagocytosis} & 1.0 & Treg effect on DC phagocytosis. \\
\bottomrule
\end{tabular}
\end{table}

\subsection{Global Weights (4+ parameters)}

\begin{table}[H]
\centering
\caption{Global weight parameters.}
\label{tab:global-weights}
\begin{tabular}{@{}llp{5.5cm}@{}}
\toprule
\textbf{Parameter} & \textbf{Default} & \textbf{Description} \\
\midrule
\texttt{w\_search\_dimension} & 0.5 & Controls adaptive search radius for agent interactions: $\mathrm{dim} = \max(5, \lfloor n_{\mathrm{tumor}} / (10 \cdot w) \rfloor)$. \\
\texttt{w\_tumour\_growth\_threshold} & 1.0 & Terminal condition threshold for mean tumour growth rate. \\
\texttt{w\_cell\_base\_death\_prob} & 0.009 & Per-step base death probability for all cells. \\
\texttt{w\_progressive\_exhaustion} & 0.05 & Rate of progressive immune cell exhaustion. \\
\texttt{w\_ici\_effectiveness} & 1.0 & ICI drug effectiveness scaling. \\
\texttt{w\_tki\_effectiveness} & 1.0 & TKI drug effectiveness scaling. \\
\texttt{w\_sex\_hormone\_cd8} & 1.0 & Sex hormone influence on CD8+ population. \\
\texttt{w\_dc\_phagocytosis\_effect} & 1.0 & DC phagocytosis effect scaling. \\
\texttt{b\_dc\_phagocytosis} & 0.0 & Bias for DC phagocytosis. \\
\texttt{w\_max\_drift} & 0.25 & Maximum sex-specific CD8+ drift over the drift period. \\
\texttt{receptor\_threshold\_variation} & 1 & Integer variation in TCR matching threshold. \\
\bottomrule
\end{tabular}
\end{table}

% ======================================================================
\section{Glucose Parameters}
\label{sec:glucose-params}

The glucose microenvironment system introduces 8 new weight parameters, all prefixed with \texttt{w\_glucose\_}. These parameters control the physical dynamics of glucose diffusion, the metabolic consumption rates, and the coupling between glucose availability and agent behavior.

\begin{longtable}{@{}llp{7cm}@{}}
\caption{Glucose weight parameters.} \label{tab:glucose-params} \\
\toprule
\textbf{Parameter} & \textbf{Default} & \textbf{Description} \\
\midrule
\endfirsthead
\toprule
\textbf{Parameter} & \textbf{Default} & \textbf{Description} \\
\midrule
\endhead
\texttt{w\_glucose\_diffusion} & 0.1 & Diffusion coefficient $D$ for the 3D Laplacian update. Stability requires $D < 1/6 \approx 0.167$; the default value of 0.1 provides moderate diffusion well within the stable regime. \\[4pt]

\texttt{w\_glucose\_decay} & 0.01 & Natural glucose decay rate per step. Applied multiplicatively as $(1 - d)$ to the entire field, modeling background metabolic consumption by tissue not explicitly represented as agents. \\[4pt]

\texttt{w\_glucose\_source\_rate} & 5.0 & Amount of glucose injected per blood vessel agent per step. Blood vessels act as point sources, creating localized glucose-rich zones that subsequently diffuse outward. \\[4pt]

\texttt{w\_glucose\_tumor\_consumption} & 2.0 & Glucose consumed per tumor cell per step. Reflects the elevated glycolytic rate characteristic of the Warburg effect~\cite{warburg1956}. Tumor cells consume approximately 4$\times$ the glucose of immune cells. \\[4pt]

\texttt{w\_glucose\_immune\_consumption} & 0.5 & Glucose consumed per immune cell per step. Represents the metabolic cost of immune cell activation and effector function. \\[4pt]

\texttt{w\_glucose\_growth\_sensitivity} & 1.0 & Sensitivity of tumor growth probability to local glucose concentration. The glucose factor is computed as $\min(1, \, C \cdot s)$, where $C$ is the local concentration and $s$ is this sensitivity parameter. \\[4pt]

\texttt{w\_glucose\_immune\_sensitivity} & 1.0 & Sensitivity of immune cell effector function to local glucose concentration. Controls how strongly glucose deprivation impairs immune cell activity. \\[4pt]

\texttt{w\_glucose\_chemotaxis\_strength} & 0.3 & Probability that an immune cell follows the glucose gradient (chemotaxis) rather than performing a random walk. At each step, with probability $p = 0.3$, the cell moves in the direction of increasing glucose concentration. \\
\bottomrule
\end{longtable}

\subsection{Parameter Interactions}

The glucose parameters interact with the existing weight parameters in several ways:

\begin{itemize}
    \item \textbf{Tumor growth modulation}: The tumor cell duplication probability (governed by \texttt{w\_tumor\_growth\_eff} and \texttt{w\_tumor\_growth\_dna}) is multiplied by a glucose factor $\min(1, C \cdot \texttt{w\_glucose\_growth\_sensitivity})$. When glucose is abundant, this factor approaches 1 and growth proceeds at the base rate; when glucose is depleted, growth is suppressed.

    \item \textbf{Immune function modulation}: Immune cell kill rates and activation probabilities can be scaled by an analogous glucose factor using \texttt{w\_glucose\_immune\_sensitivity}, coupling metabolic resource availability to immune effectiveness.

    \item \textbf{Spatial competition}: The balance between \texttt{w\_glucose\_source\_rate} (supply from blood vessels) and \texttt{w\_glucose\_tumor\_consumption} (demand from tumor cells) determines the spatial glucose profile. High tumor consumption relative to supply creates depletion zones that impair both further tumor growth and immune infiltration.
\end{itemize}

% ======================================================================
\section{Bayesian Optimization}
\label{sec:bayesian-optimization}

\subsection{Framework}

The model uses Optuna~\cite{optuna} for Bayesian hyperparameter optimization. Optuna's Tree-structured Parzen Estimator (TPE) sampler efficiently explores high-dimensional parameter spaces by building probabilistic models of the objective function landscape.

\subsection{Objective Function}

The optimization objective is to minimize the prediction error of the model against the ARON-1 clinical dataset. For each Optuna trial, the following process is executed:

\begin{enumerate}
    \item A set of weight parameters is suggested by the TPE sampler within predefined bounds.
    \item For each patient case in the ARON-1 dataset, a simulation is run with the patient's clinical features (sex, BMI, treatment, immune cell concentrations) and the suggested weight parameters.
    \item The simulation outcome (survival/death, number of steps) is compared against the clinical outcome.
    \item The per-case error is computed as:
    \begin{equation}
        \text{error}_i = \begin{cases}
            (\text{steps}_i - \text{OS}_i)^2 & \text{if prediction matches outcome} \\
            \text{HIGH\_PENALTY}^2 & \text{if prediction contradicts outcome}
        \end{cases}
    \end{equation}
    where $\text{OS}_i$ is the observed overall survival (in steps) and $\text{HIGH\_PENALTY} = 500$.
    \item The trial objective value is the mean error across all cases.
\end{enumerate}

\subsection{Search Space}

The search space is defined in the \texttt{suggest\_parameters()} function and covers the full 92-parameter weight space:

\begin{itemize}
    \item \textbf{Positive weights} (62 parameters): Sampled log-uniformly in $[0.25, 4.0]$. The log-uniform distribution ensures adequate exploration of both small and large scaling factors.
    \item \textbf{Biases} (15 parameters): Sampled uniformly in $[-1.0, 1.0]$.
    \item \textbf{Effectiveness weights}: \texttt{w\_ici\_effectiveness} and \texttt{w\_tki\_effectiveness} in $[0.0, 1.0]$.
    \item \textbf{Threshold weights}: \texttt{w\_tumour\_growth\_threshold} in $[0.75, 0.99]$; \texttt{w\_cell\_base\_death\_prob} in $[0.001, 0.1]$; \texttt{w\_progressive\_exhaustion} in $[0.005, 0.3]$.
    \item \textbf{Integer parameters}: \texttt{receptor\_threshold\_variation} in $\{0, 1, 2, 3\}$.
    \item \textbf{Glucose parameters} (8 parameters):
    \begin{itemize}
        \item \texttt{w\_glucose\_diffusion} in $[0.01, 0.16]$ (respecting the $D < 1/6$ stability constraint).
        \item \texttt{w\_glucose\_decay} in $[0.001, 0.1]$.
        \item \texttt{w\_glucose\_source\_rate} in $[1.0, 20.0]$.
        \item \texttt{w\_glucose\_tumor\_consumption} in $[0.5, 10.0]$.
        \item \texttt{w\_glucose\_immune\_consumption} in $[0.1, 2.0]$.
        \item \texttt{w\_glucose\_growth\_sensitivity} in $[0.1, 5.0]$.
        \item \texttt{w\_glucose\_immune\_sensitivity} in $[0.1, 5.0]$.
        \item \texttt{w\_glucose\_chemotaxis\_strength} in $[0.05, 0.8]$.
    \end{itemize}
\end{itemize}

\subsection{Optimization Configuration}

The optimization runs for $N = 250$ trials, each evaluating every patient case in the ARON-1 dataset. Training uses \texttt{MPI.COMM\_SELF} to create isolated single-rank model instances, avoiding MPI communication overhead during optimization. Trial errors and the best-so-far error are logged to \texttt{data/trial\_errors.csv} for convergence analysis.

\begin{lstlisting}[caption={Optuna study creation and execution.}, label={lst:optuna-study}]
study = optuna.create_study(direction="minimize")
study.optimize(wrapped_objective, n_trials=250)
\end{lstlisting}

% ======================================================================
\section{Clinical Dataset}
\label{sec:clinical-dataset}

\subsection{ARON-1 Study}

The ARON-1 study is a clinical trial of advanced renal cell carcinoma patients treated with immune checkpoint inhibitors (ICI), tyrosine kinase inhibitors (TKI), or their combination~\cite{rini2019, motzer2019}. The dataset provides the ground truth against which the model is calibrated.

\subsection{Dataset Structure}

The clinical dataset is stored in \texttt{src/learning/data/ARON.csv} and contains the following features for each patient case:

\begin{itemize}
    \item \textbf{Patient features}: sex (\texttt{'F'}/\texttt{'M'}), BMI, treatment type (ICI, TKI, or ICI+TKI).
    \item \textbf{Immune cell concentrations}: 13 cell type concentrations matching the \texttt{PatientParameters} fields.
    \item \textbf{Outcome}: Overall survival (OS) in time units, and a binary death indicator.
\end{itemize}

The dataset is loaded with Pandas, with the treatment column parsed as a list of strings to support combination therapies:

\begin{lstlisting}[caption={Dataset loading.}]
def load_dataset(file_path: str) -> pd.DataFrame:
    def treatment_to_list(treatment_str):
        return [t.strip() for t in treatment_str.split('+')]
    return pd.read_csv(file_path,
                       converters={"treatment": treatment_to_list})
\end{lstlisting}

Each row of the dataset is used to parameterize a simulation run during optimization, with the patient's clinical features passed as keyword arguments to \texttt{RCCModel}. The simulation's survival prediction is then compared against the clinical outcome to compute the trial error.
